\SourceLeaf{286}%
\noindent wobei $F$ den Flächeninhalt der ganzen Figur und $S$ den der gesamten Fläche bezeichnet,
\[
= (2\chi - \sigma) \, . \, 360^0 - \frac{F}{S} \,.\,720^0.
\]
Diese Regel schließt die oben angegebene ein, da das unbeschränkte Gebiet in der Ebene den Charakter 0 hat und damit $\frac{F}{S} = 0$ wird. Oben wurden die Vorzeichen dort gezählt, wo man die Schnittpunkte zum ersten Mal traf. Dort sind sie den hier verwendeten entgegengesetzt.

\begin{center}
[7.]
\end{center}

\begin{center}\textbf{EINE ANDERE ART, DIESE GEGENSTÄNDE ZU BEHANDELN.}\end{center}

Es sei $S$ die ganze Kurve und $s$ ein unbestimmtes Stück. Im Anfangspunkt ziehe man eine Tangente und bezeichne mit $(0, s)$ den Winkel, den die Gerade vom Anfangspunkt zum Endpunkt des Stücks $s$ bildet. Offensichtlich ändert sich $(0, s)$ von $s = 0$ an stetig, solange die Kurve nicht wieder zum Anfangspunkt zurückgekehrt ist. Wir nehmen an, dass dies erst geschieht, wenn $s = S$ geworden ist, und der letzte Wert, bevor $s$ diesen Wert erreicht, ist
\[
(0, s) = \pm 180^0,
\]
je nachdem, ob man die Winkel zur inneren oder äußeren Seite der Linie hin wachsen lässt.

Man halte $s$ konstant, ziehe von jedem Punkt der Linie, die als Endpunkt des kleineren Stücks $s'$ betrachtet wird, eine Gerade zum Endpunkt des Stücks $s$ und bezeichne deren Neigung zur ersten Tangente mit $(s', s)$. Solange der Endpunkt von $s'$ nicht mit dem von $s$ zusammenfällt, ändert sich bei stetiger Änderung von $s'$ auch $(s', s)$ stetig, sofern sich das Stück $s$ nicht selbst schneidet, bis $s' = s$. Für diesen Grenzwert selbst ist $(s', s)$ genau genommen unbestimmt. Solange aber $s' < s$ gilt, ist der Wert bestimmt, so nahe die Werte auch beieinanderliegen. Je näher $s'$ an $s$ herankommt, desto mehr nähert sich dieser Wert der Neigung der Tangente zur Tangente im Anfangspunkt. Nun sei $s = S - \omega$, und $s'$ ändere sich von $0$ bis $S - \omega - \omega'$. Offenbar kommt die Änderung von $s'$ dadurch $180^0$ umso näher, je kleiner $\omega, \omega'$ sind, sodass $(S - \omega - \omega', S - \omega)$ dem
